% Generated draft: requires mathematical and source review.
\documentclass[11pt]{article}
\usepackage[margin=25mm]{geometry}
\usepackage{fontspec}
\setmainfont{Latin Modern Roman}
\usepackage{amsmath,amssymb,mathtools,unicode-math}
\setmathfont{Latin Modern Math}
\usepackage{enumitem,needspace,xurl,hyperref}
\setlength{\parindent}{0pt}
\setlength{\parskip}{5pt}
\setlength{\emergencystretch}{3em}
\newcommand{\N}{\mathbb{N}}
\newcommand{\Z}{\mathbb{Z}}
\newcommand{\Q}{\mathbb{Q}}
\newcommand{\R}{\mathbb{R}}
\newcommand{\C}{\mathbb{C}}
\newcommand{\F}{\mathbb{F}}
\newcommand{\A}{\mathbb{A}}
\newcommand{\PP}{\mathbb{P}}
\newcommand{\E}{\mathbb{E}}
\newcommand{\eps}{\varepsilon}
\newcommand{\dd}{\,\mathrm{d}}
\begin{document}
\section*{UnsolvedMath ID: 20002250}
\textbf{Source problem code:} catalogue entry
\subsection{Definitions and mathematical statement}
This extraction merges Questions 10--14 of the 2014 workshop notes; the initial ``210'' is not a new problem number. The supplied aid's proved reduction concerns only the literal dual-covering formulation in Question 13.
\begin{enumerate}
\item Question 10 asks for the consistency strength of simultaneous tree property at \(\aleph_2\) and \(\aleph_3\). At a regular uncountable \(\kappa\), the tree property says every height-\(\kappa\) tree with levels of size less than \(\kappa\) has a cofinal branch. The historical note gives a weakly compact cardinal above a supercompact as an upper bound, attributed to Abraham, and a Woodin-cardinal lower comparison via Foreman--Magidor--Schindler. The parenthetical ``2 is open'' is not sufficiently explained to interpret further.
\item Question 11 asks whether there is a unique \(L(\R,\mu)\) internally satisfying that \(\mu\) is a normal fine measure on \(P_{\omega_1}(\R)\), and asks its consistency strength. Fineness means each real belongs to \(\mu\)-almost every countable set of reals; normality makes regressive selections constant almost everywhere, with internal-model conventions required. It reports comparability: for two such models, \(\mathcal P(\R)\cap L(\R,\mu)\subseteq L(\R,\nu)\), or conversely. This is not stated as equality or uniqueness. The extracted lower-bound string ``\(\omega2\) Woodins'' does not reliably distinguish a superscript from a subscript.
\item Question 12 asks whether bounded Martin's Maximum (BMM) implies existence of the mouse printed ``0¶''. It reports Schindler's implication from BMM to an inner model with a strong cardinal and Woodin's consistency construction from \(\omega+1\) Woodin cardinals plus a measurable. The source labels these ``upper'' and ``lower'' respectively; the precise relative-consistency directions and the mouse notation need verification rather than automatic relabelling.
\item Question 13 asks whether large cardinals can yield failure of a proposed dual-covering principle at \((M,\aleph_3,\aleph_1)\), for suitable inner models. The printed examples refer to one-Woodin \(K\) under absence of a proper-class model with a Woodin, and again under absence of one with a strong cardinal. These descriptions are not precise enough to identify the intended core-model hypotheses, and are not repaired by invention.
\item Question 14 concerns strong condensation: for each cardinal \(\kappa>\omega\), a bijection \(h:\kappa\to H(\kappa)\) has the displayed property
\[
\pi[X\cap h]=h\upharpoonright\operatorname{ot}(X\cap\kappa),\qquad
X\prec(H(\kappa),h).
\]
Here \(H(\kappa)\) consists of sets of hereditary size less than \(\kappa\). The image notation is naturally read in the collapse direction, whereas the raw text calls \(\pi\) the uncollapse; that typing convention remains unresolved. If an inner model \(N\) satisfies strong condensation and covering fails relative to \(N\), must \(N^\#\) exist? The sharp is restored by the supplied verification; the exact covering principle and sharp framework are not defined in the extract.
\end{enumerate}

For Question 13, make the literal quantifiers explicit. Let \(M\) be a transitive inner model of ZFC containing all ordinals, let \(\lambda>0\) be an ordinal, \(\delta\geq1\) a cardinal, and \(f:\lambda^{<\omega}\to\lambda\) any ambient function, not necessarily in \(M\). Call a \emph{set} of ordinals \(X\) \(f\)-closed if
\[
f``(X\cap\lambda)^{<\omega}\subseteq X.
\]
Let \(\operatorname{DCov}_f(M,\lambda,\delta)\) say every such \(X\) is a union of at most \(\delta\) sets belonging to \(M\), with the family indexed in the ambient universe. The supplied aid reports the equivalence
\[
\operatorname{DCov}_f(M,\lambda,\delta)
\quad\Longleftrightarrow\quad
\forall A\subseteq\operatorname{Ord}\text{ a set},\quad
A\text{ is a union of at most }\delta\text{ members of }M.
\]
For the nontrivial direction, encode arbitrary \(A\) as \(X=\lambda\cup\{\lambda+\alpha:\alpha\in A\}\). This \(X\) is \(f\)-closed regardless of \(f\); from each \(M\)-piece of \(X\), separation and ordinal-addition inversion recover an \(M\)-piece of \(A\). Thus the literal principle is independent of the chosen \(f\) and of nonzero \(\lambda\). Allowing empty pieces removes the distinction between at most \(\delta\) and exactly \(\delta\) pieces.

This quantifier collapse is not an answer to the intended large-cardinal consistency problem. Requiring \(X\subseteq\lambda\), bounding its size, using a different closure domain, or demanding an internal indexing family changes the assertion. No such repair is silently supplied, and no current resolution of Questions 10--12 or 14 is claimed.
\subsection{Short English statement}
The record concatenates five inner-model-theory questions with several damaged conventions. The supplied aid gives one precise result: under the literal printed quantifiers, dual covering reduces to uniform coverability of every set of ordinals and no longer depends on its closure function or nonzero domain. That does not settle a corrected consistency problem.
\subsection{Sources}
\begin{itemize}
\item \emph{Problem Session Notes}, AIM workshop \emph{Descriptive Inner Model Theory}, 2--6 June 2014, Questions 10--14: \url{https://aimath.org/pastworkshops/innermodelproblems.pdf}.
\item Named attributions in the notes: Abraham; Foreman--Magidor--Schindler; Schindler; Woodin. No full theorem references are supplied. The historical bounds and comparability statement are retained as source reports.
\item Local registry \texttt{data/unsolvedmath/problems.json}, permanent ID 20002250, code \texttt{catalogue entry}; canonical \texttt{aim-logic-notes.json}, \texttt{notes[25]}; updated 14 August 2026. Clean statement null; statement status reconstructed unverified. The supplied aid restores the restriction symbol and \(N^\#\) while preserving the original damaged record separately.
\item Attempt 2: partial result, partially solved status, rigor label proved, novelty confidence low. The machine-generated aid is not peer reviewed and requires independent verification before citation; the original literature field is empty.
\end{itemize}
\end{document}
